% =============================================================================
% PARTE III: TECNOLOGÍAS DE DETECCIÓN AUTOMÁTICA Y SISTEMAS DE EXTINCIÓN
% =============================================================================

% --- SLIDE 29: ARQUITECTURA DE DETECCIÓN ---
\begin{frame}{Topologías de Detección: Convencional frente a Direccionable}
  \begin{columns}[T]
    \begin{column}{0.48\textwidth}
      \begin{block}{\textbf{Sistema Convencional (Zonas Radiales)}}
        \scriptsize
        $\bullet$ Circuito radial de 2 hilos con resistencia fin de línea ($R_{\mathrm{EOL}}$).\\
        $\bullet$ Alarma por caída de impedancia en toda la zona.\\
        $\bullet$ \textbf{Limitación:} No individualiza el detector activado; requiere búsqueda visual en el local.
      \end{block}
      \vspace{0.15cm}
      \centering
      \resizebox{0.95\columnwidth}{!}{%
        \begin{tikzpicture}[font=\tiny]
          % Central
          \draw[fill=utnblue!15, draw=utnblue, line width=1.2pt, rounded corners=2pt] (0,-0.7) rectangle (1.2,0.7) node[midway, align=center, font=\fontsize{5}{6}\selectfont\bfseries, utnblue] {Central\\Convencional};
          % Línea radial
          \draw[line width=1.2pt, darkcharcoal] (1.2,0.3) -- (4.2,0.3);
          \draw[line width=1.2pt, darkcharcoal] (1.2,-0.3) -- (4.2,-0.3);
          % Detectores
          \draw[fill=white, draw=darkcharcoal, line width=1pt] (2.0,0) circle (0.28) node[font=\fontsize{4}{5}\selectfont\bfseries] {D1};
          \draw[line width=0.8pt] (2.0,0.3) -- (2.0,0.28);
          \draw[line width=0.8pt] (2.0,-0.3) -- (2.0,-0.28);
          
          \draw[fill=firered!20, draw=firered, line width=1.2pt] (3.1,0) circle (0.28) node[font=\fontsize{4}{5}\selectfont\bfseries, firered] {D2};
          \draw[line width=0.8pt] (3.1,0.3) -- (3.1,0.28);
          \draw[line width=0.8pt] (3.1,-0.3) -- (3.1,-0.28);
          
          % Resistencia EOL
          \draw[line width=0.8pt] (4.2,0.3) -- (4.2,0.18);
          \draw[line width=0.8pt] (4.2,-0.18) -- (4.2,-0.3);
          \draw[fill=amberyellow!30, draw=darkcharcoal, line width=0.8pt] (3.85,-0.18) rectangle (4.55,0.18) node[midway, font=\fontsize{4}{5}\selectfont\bfseries] {$R_{\mathrm{EOL}}$};
          
          % Alerta
          \node[draw=firered, fill=firered!10, rounded corners=2pt, font=\fontsize{4.5}{5.5}\selectfont, align=center, text=firered] at (2.6,-0.9) {Disparo en D2 $\rightarrow$ Solo indica ``Zona 1''};
        \end{tikzpicture}
      }
    \end{column}

    \begin{column}{0.48\textwidth}
      \begin{block}{\textbf{Sistema Analógico-Direccionable (Anillo)}}
        \scriptsize
        $\bullet$ Lazo cerrado con retorno a central.\\
        $\bullet$ Cada sensor tiene microprocesador e ID.\\
        $\bullet$ Supervisión analógica continua.\\
        $\bullet$ \textbf{Ventaja:} Localización puntual y tolerancia a corte.
      \end{block}
      \vspace{0.15cm}
      \centering
      \resizebox{0.95\columnwidth}{!}{%
        \begin{tikzpicture}[font=\tiny]
          % Central FACP
          \draw[fill=successgreen!15, draw=successgreen!80!black, line width=1.2pt, rounded corners=2pt] (0,-0.7) rectangle (1.2,0.7) node[midway, align=center, font=\fontsize{5}{6}\selectfont\bfseries, successgreen!80!black] {Central\\Direccionable};
          % Lazo en anillo Clase A
          \draw[line width=1.2pt, successgreen!70!black] (1.2,0.4) -- (2.0,0.6) -- (3.6,0.6) -- (4.2,0) -- (3.6,-0.6) -- (2.0,-0.6) -- (1.2,-0.4);
          % Sensores con dirección
          \draw[fill=white, draw=successgreen!80!black, line width=1pt] (2.2,0.6) circle (0.24) node[font=\fontsize{4}{5}\selectfont\bfseries] {\#01};
          \draw[fill=white, draw=successgreen!80!black, line width=1pt] (3.4,0.6) circle (0.24) node[font=\fontsize{4}{5}\selectfont\bfseries] {\#02};
          \draw[fill=white, draw=successgreen!80!black, line width=1pt] (4.2,0) circle (0.24) node[font=\fontsize{4}{5}\selectfont\bfseries] {\#03};
          \draw[fill=white, draw=successgreen!80!black, line width=1pt] (3.4,-0.6) circle (0.24) node[font=\fontsize{4}{5}\selectfont\bfseries] {\#04};
          \draw[fill=white, draw=successgreen!80!black, line width=1pt] (2.2,-0.6) circle (0.24) node[font=\fontsize{4}{5}\selectfont\bfseries] {\#05};
          
          % Flechas bidireccionales
          \draw[->, >=stealth, line width=0.8pt, successgreen!80!black] (1.5,0.48) -- (1.8,0.55);
          \draw[<-, >=stealth, line width=0.8pt, successgreen!80!black] (1.5,-0.48) -- (1.8,-0.55);
          
          % Texto centro
          \node[draw=successgreen!80!black, fill=successgreen!10, rounded corners=2pt, font=\fontsize{4.5}{5.5}\selectfont, align=center, text=darkcharcoal] at (2.6,-1.2) {Reporta ID exacto y nivel analógico continuo};
        \end{tikzpicture}
      }
    \end{column}
  \end{columns}
\end{frame}

% --- SLIDE 30: DETECTOR ÓPTICO ---
\begin{frame}{Detector Óptico Puntual de Humo (Efecto Tyndall)}
  \begin{columns}[c]
    \begin{column}{0.50\textwidth}
      \small
      Principio físico de dispersión de la luz en cámara oscura:
      \begin{itemize}
        \scriptsize
        \item Cámara laberíntica opaca que bloquea la luz ambiental exterior pero permite la libre entrada de humo.
        \item Emisor LED infrarrojo y fotoreceptor orientados perpendicularmente.
        \item \textbf{En reposo (aire limpio):} El haz IR viaja en línea recta y es absorbido por la trampa de luz.
        \item \textbf{En alarma (humo presente):} Las partículas de humo dispersan fotones hacia el fotodiodo receptor.
        \item \textbf{Aplicación óptima:} Fuegos lentos o latentes con partículas visibles.
      \end{itemize}
    \end{column}

    \begin{column}{0.48\textwidth}
      \centering
      \resizebox{0.92\columnwidth}{!}{%
        \begin{tikzpicture}[scale=0.95]
          % Carcasa exterior circular
          \draw[line width=1.5pt, darkcharcoal, fill=lightbg] (0,0) circle (2.2);
          
          % Laberinto óptico (aletas deflectoras periféricas)
          \foreach \ang in {15, 35, 55, 75, 105, 125, 145, 165, 195, 215, 235, 255, 285, 305, 325, 345} {
            \draw[line width=1.8pt, darkcharcoal!80] (\ang:2.2) -- (\ang-15:1.7);
          }
          \draw[line width=0.8pt, darkcharcoal!40, dashed] (0,0) circle (1.7);
          
          % Emisor LED IR
          \draw[fill=utnblue, draw=darkcharcoal, line width=1pt, rounded corners=1pt] (-1.6,-0.3) rectangle (-0.9,0.3);
          \node[font=\tiny\bfseries, text=white] at (-1.25,0) {LED IR};
          
          % Trampa de luz
          \draw[fill=darkcharcoal, draw=darkcharcoal, line width=1pt] (1.3,-0.35) -- (1.6,-0.2) -- (1.3,0) -- (1.6,0.2) -- (1.3,0.35) -- (1.1,0.35) -- (1.1,-0.35) -- cycle;
          \node[font=\fontsize{4}{5}\selectfont\bfseries, text=darkcharcoal, right] at (1.1,0.5) {Trampa};
          
          % Haz directo (rojo)
          \draw[->, >=stealth, line width=1.5pt, red!80] (-0.9,0) -- (1.1,0);
          
          % Receptor Fotodiodo
          \draw[fill=fireorange, draw=darkcharcoal, line width=1pt, rounded corners=1pt] (-0.6,-1.6) rectangle (0.6,-1.0);
          \node[font=\fontsize{4.5}{5.5}\selectfont\bfseries, text=white] at (0,-1.3) {Fotodiodo};
          
          % Nube de humo en el centro
          \fill[gray!40, opacity=0.7] (0,0) circle (0.45);
          \foreach \x/\y in {-0.2/0.1, 0.1/0.15, -0.1/-0.2, 0.2/-0.1, 0/0} {
            \fill[darkcharcoal!70] (\x,\y) circle (0.04);
          }
          \node[font=\fontsize{5}{6}\selectfont\bfseries, darkcharcoal, above] at (0,0.45) {Humo};
          
          % Rayos dispersados Tyndall a 90 grados
          \draw[->, >=stealth, dashed, line width=1.2pt, fireorange] (0.05,-0.05) -- (0,-0.95);
          \draw[->, >=stealth, dashed, line width=1pt, fireorange] (-0.15,-0.1) -- (-0.15,-0.95);
          \draw[->, >=stealth, dashed, line width=1pt, fireorange] (0.2,-0.1) -- (0.15,-0.95);
        \end{tikzpicture}
      }
    \end{column}
  \end{columns}
\end{frame}

% --- SLIDE 31: DETECTOR IÓNICO ---
\begin{frame}{Detector Iónico de Humo por Americio-241 ($^{241}\mathrm{Am}$)}
  \begin{columns}[c]
    \begin{column}{0.50\textwidth}
      \small
      Principio de conducción eléctrica en cámara ionizada:
      \begin{itemize}
        \scriptsize
        \item Fuente radiactiva encapsulada de $^{241}\mathrm{Am}$ que emite partículas alfa ($\alpha$).
        \item Las partículas $\alpha$ ionizan las moléculas de aire ($N_2, O_2$), separándolas en iones positivos ($+$) y electrones libres ($-$).
        \item Se genera una \textbf{corriente de reposo} ($I_{\mathrm{reposo}} \sim \mathrm{pA}$) entre las dos placas a potencial fijo.
        \item Las partículas de humo capturan los iones móviles lo que produce una caída abrupta de corriente.
        \item \textbf{Especialidad:} Fuegos abiertos de combustión rápida con llama viva.
      \end{itemize}
    \end{column}

    \begin{column}{0.48\textwidth}
      \centering
      \resizebox{0.92\columnwidth}{!}{%
        \begin{tikzpicture}[scale=0.9]
          % Cámara exterior mallada
          \draw[line width=1.5pt, darkcharcoal, fill=lightbg] (0,0) rectangle (4.2,3.2);
          \draw[dashed, line width=0.8pt, darkcharcoal!50] (0,0.8) -- (0,2.4);
          \draw[dashed, line width=0.8pt, darkcharcoal!50] (4.2,0.8) -- (4.2,2.4);
          \node[font=\fontsize{4.5}{5.5}\selectfont, darkcharcoal, above] at (2.1,3.2) {Cámara de Ionización};

          % Placa Superior (+V)
          \draw[fill=firered!80, draw=darkcharcoal, line width=1pt] (0.6,2.7) rectangle (3.6,2.95);
          \node[font=\tiny\bfseries, text=white] at (2.1,2.82) {Electrodo Colector ($+V$)};

          % Placa Inferior (Tierra / Circuito sensor)
          \draw[fill=utnblue!80, draw=darkcharcoal, line width=1pt] (0.6,0.25) rectangle (3.6,0.5);
          \node[font=\tiny\bfseries, text=white] at (2.1,0.37) {Electrodo Base (Sensor $I$)};

          % Fuente Americio 241
          \draw[fill=amberyellow, draw=darkcharcoal, line width=1pt] (1.8,0.5) rectangle (2.4,0.75);
          \node[font=\fontsize{4}{5}\selectfont\bfseries, text=darkcharcoal] at (2.1,0.62) {$^{241}\mathrm{Am}$};

          % Partículas alfa y pares de iones
          \draw[->, >=stealth, line width=1pt, amberyellow!90!black] (2.0,0.75) -- (1.2,1.8) node[midway, left, font=\fontsize{5}{6}\selectfont\bfseries] {$\alpha$};
          \draw[->, >=stealth, line width=1pt, amberyellow!90!black] (2.2,0.75) -- (3.0,1.8) node[midway, right, font=\fontsize{5}{6}\selectfont\bfseries] {$\alpha$};

          % Iones libres
          \node[circle, fill=utncyan, text=white, inner sep=1pt, font=\tiny\bfseries] at (1.3,1.4) {$+$};
          \node[circle, fill=fireorange, text=white, inner sep=1pt, font=\tiny\bfseries] at (1.6,2.2) {$-$};
          \node[circle, fill=utncyan, text=white, inner sep=1pt, font=\tiny\bfseries] at (2.9,1.3) {$+$};
          \node[circle, fill=fireorange, text=white, inner sep=1pt, font=\tiny\bfseries] at (2.7,2.3) {$-$};

          % Humo entrando
          \draw[->, >=stealth, line width=1.2pt, darkcharcoal!70] (0.0,1.6) -- (0.8,1.6) node[midway, above, font=\fontsize{4}{5}\selectfont\bfseries] {Humo};
          \draw[->, >=stealth, line width=1.2pt, darkcharcoal!70] (4.2,1.6) -- (3.4,1.6) node[midway, above, font=\fontsize{4}{5}\selectfont\bfseries] {Humo};
          \fill[gray!50, opacity=0.7] (2.1,1.7) circle (0.35);
          \node[font=\fontsize{4.5}{5.5}\selectfont\bfseries, darkcharcoal] at (2.1,1.7) {Captura};

          % Medidor de corriente
          \draw[fill=white, draw=darkcharcoal, line width=0.8pt] (3.6,0.37) -| (4.0,-0.3) -- (3.5,-0.3);
          \draw[fill=white, draw=darkcharcoal, line width=0.8pt] (2.1,0.25) |- (2.7,-0.3);
          \node[draw, circle, fill=cardbg, line width=0.8pt, inner sep=2pt, font=\fontsize{4}{5}\selectfont\bfseries] at (3.1,-0.3) {$\mathrm{pA}$};
          \node[font=\fontsize{4}{5}\selectfont, firered, below] at (3.1,-0.55) {$\Delta I \downarrow \rightarrow \text{Alarma}$};
        \end{tikzpicture}
      }
    \end{column}
  \end{columns}
\end{frame}

% --- SLIDE 32: BARRERA ÓPTICA LINEAL ---
\begin{frame}{Detección Lineal por Haz Proyectado (Barreras Ópticas)}
  \begin{columns}[c]
    \begin{column}{0.48\textwidth}
      \small
      Tecnología orientada a locales industriales y techos de gran altura:
      \begin{itemize}
        \scriptsize
        \item Compuesto por un emisor de haz infrarrojo y un receptor fotosensible.
        \item Cubre distancias longitudinales de hasta $\mathbf{100\,\m}$.
        \item \textbf{Principio de funcionamiento:} Atenuación del haz de luz por las particulas de humo que la absorben.
        \item \textbf{Ventaja operativa:} Inmune al retardo por estratificación térmica.
      \end{itemize}
    \end{column}

    \begin{column}{0.50\textwidth}
      \centering
      \resizebox{0.95\columnwidth}{!}{%
        \begin{tikzpicture}[scale=0.9]
          % Silueta nave industrial con techo a dos aguas
          \draw[line width=1.2pt, darkcharcoal, fill=lightbg] (0,0) -- (0,2.5) -- (2.5,3.3) -- (5.0,2.5) -- (5.0,0) -- cycle;
          % Suelo
          \draw[line width=1.5pt, darkcharcoal!70] (-0.2,0) -- (5.2,0);
          \node[font=\fontsize{4.5}{5.5}\selectfont, darkcharcoal] at (2.5,0.2) {Piso de Planta};

          % Capa de estratificación de humo bajo techo
          \fill[gray!40, opacity=0.85] (0.2,1.8) .. controls (1.5,2.4) and (3.5,2.4) .. (4.8,1.8) -- (4.8,2.4) -- (2.5,3.1) -- (0.2,2.4) -- cycle;
          \node[font=\fontsize{5}{6}\selectfont\bfseries, darkcharcoal] at (2.5,2.6) {Capa de Humo Estratificado};

          % Columna térmica ascendente
          \draw[->, >=stealth, line width=1pt, firered, dashed] (2.5,0.4) .. controls (2.3,1.0) .. (2.5,1.8);
          \draw[->, >=stealth, line width=1pt, firered, dashed] (2.3,0.4) .. controls (2.1,1.0) .. (2.3,1.8);
          \draw[->, >=stealth, line width=1pt, firered, dashed] (2.7,0.4) .. controls (2.9,1.0) .. (2.7,1.8);
          \node[font=\fontsize{4.5}{5.5}\selectfont\bfseries, firered] at (2.5,0.6) {Foco Térmico};

          % Emisor IR en pared izquierda
          \draw[fill=utnblue, draw=darkcharcoal, line width=0.8pt, rounded corners=1pt] (0.15,1.9) rectangle (0.75,2.3);
          \node[font=\fontsize{4.5}{5.5}\selectfont\bfseries, text=white] at (0.45,2.1) {Emisor};

          % Receptor / Prisma en pared derecha
          \draw[fill=utnblue, draw=darkcharcoal, line width=0.8pt, rounded corners=1pt] (4.25,1.9) rectangle (4.85,2.3);
          \node[font=\fontsize{4.5}{5.5}\selectfont\bfseries, text=white] at (4.55,2.1) {Receptor};

          % Haz infrarrojo colimado
          \draw[line width=1.8pt, firered, dashed] (0.75,2.1) -- (4.25,2.1);
          \node[font=\fontsize{4.5}{5.5}\selectfont\bfseries, firered, above] at (2.5,2.1) {Haz IR};

          % Cota de distancia
          \draw[<->, >=stealth, line width=0.8pt, darkcharcoal] (0.75,1.4) -- (4.25,1.4) node[midway, fill=lightbg, font=\fontsize{4.5}{5.5}\selectfont\bfseries] {Alcance hasta $100\,\m$};
        \end{tikzpicture}
      }
    \end{column}
  \end{columns}
\end{frame}

% --- SLIDE 33: ASPIRACIÓN ASD / VESDA ---
\begin{frame}{Detección por Aspiración de Alta Sensibilidad (ASD / VESDA)}
  \begin{columns}[c]
    \begin{column}{0.50\textwidth}
      \small
      El estándar de detección de alta fidelidad en infraestructuras críticas:
      \begin{itemize}
        \scriptsize
        \item Red de tuberías con orificios a lo largo de todo el local.
        \item Aspirador de alta eficiencia que succiona aire continuo hacia el módulo de detección central.
        \item Filtro de doble etapa: separa partículas gruesas y polvo ambiental.
        \item \textbf{Cámara láser de alta precisión:} Detección por dispersión de partículas con sensibilidad de hasta $\mathbf{0.005\%\,\mathrm{obs/m}}$.
        \item Alerta hasta 60 minutos antes del conato con llama.
      \end{itemize}
    \end{column}

    \begin{column}{0.48\textwidth}
      \centering
      \resizebox{0.92\columnwidth}{!}{%
        \begin{tikzpicture}[scale=0.85]
          % Módulo detector VESDA
          \draw[line width=1.5pt, fill=cardbg, draw=utnblue, rounded corners=4pt] (0,0) rectangle (4.2,4.0);
          \node[fill=utnblue, text=white, rounded corners=2pt, font=\fontsize{5.5}{6.5}\selectfont\bfseries, minimum width=3.8cm, minimum height=0.45cm] at (2.1,3.7) {MÓDULO CENTRAL VESDA};

          % Cámara láser
          \draw[fill=firered!15, draw=firered, line width=1pt, rounded corners=2pt] (0.4,2.4) rectangle (3.8,3.2);
          \node[font=\fontsize{5}{6}\selectfont\bfseries, firered] at (2.1,2.95) {Cámara de Análisis Láser};
          \node[font=\fontsize{4}{5}\selectfont, darkcharcoal] at (2.1,2.6) {Sensibilidad $0.005\%\,\mathrm{obs/m}$};

          % Filtro de partículas
          \draw[fill=amberyellow!20, draw=amberyellow!80!black, line width=1pt, rounded corners=2pt] (0.4,1.4) rectangle (3.8,2.1);
          \node[font=\fontsize{5}{6}\selectfont\bfseries, amberyellow!80!black] at (2.1,1.85) {Filtro de Doble Etapa};
          \node[font=\fontsize{4}{5}\selectfont, darkcharcoal] at (2.1,1.55) {Eliminación de polvo y pelusa};

          % Aspirador de turbina
          \draw[fill=utncyan!20, draw=utncyan, line width=1pt, rounded corners=2pt] (0.4,0.4) rectangle (3.8,1.1);
          \node[font=\fontsize{5}{6}\selectfont\bfseries, utncyan] at (2.1,0.85) {Aspirador de Turbina Continua};
          \node[font=\fontsize{4}{5}\selectfont, darkcharcoal] at (2.1,0.58) {Presión negativa de muestreo};

          % Flujo interno
          \draw[->, >=stealth, line width=1.5pt, utnblue] (2.1,1.1) -- (2.1,1.4);
          \draw[->, >=stealth, line width=1.5pt, utnblue] (2.1,2.1) -- (2.1,2.4);

          % Tubería externa de entrada
          \draw[line width=2.5pt, firered] (0.9,-0.6) -- (0.9,0.4);
          \draw[->, >=stealth, line width=1.5pt, firered] (0.9,-0.6) -- (0.9,0.1);
          \node[font=\fontsize{4.5}{5.5}\selectfont\bfseries, firered, left] at (0.8,-0.3) {Red de Tuberías};

          % Escape de aire
          \draw[line width=2pt, darkcharcoal!50] (3.3,3.2) -- (3.3,4.4);
          \draw[->, >=stealth, line width=1.2pt, darkcharcoal!60] (3.3,3.2) -- (3.3,4.3);
          \node[font=\fontsize{4.5}{5.5}\selectfont, darkcharcoal, right] at (3.4,4.2) {Escape};
        \end{tikzpicture}
      }
    \end{column}
  \end{columns}
\end{frame}

% --- SLIDE 34: TÉRMICA, CABLE SENSOR Y LLAMA ---
\begin{frame}{Detección Térmica, Cable Sensor Lineal y Detectores de Llama}
  \begin{columns}[T]
    \begin{column}{0.48\textwidth}
      \begin{block}{\textbf{Cable Sensor Lineal Térmico}}
        \scriptsize
        $\bullet$ Dos conductores elásticos bajo torsión mecánica, aislados por polímero termosensible.\\
        $\bullet$ Al superar la temperatura crítica ($68\text{--}105\,\Celsius$), el polímero funde $\rightarrow$ \textbf{cortocircuito calibrado}.\\
        $\bullet$ La central mide la resistencia del lazo y calcula la \textbf{distancia en metros} del siniestro.\\
        $\bullet$ \textbf{Uso:} Bandejas de cables y túneles.
      \end{block}
      \vspace{0.1cm}
      \centering
      \resizebox{0.95\columnwidth}{!}{%
        \begin{tikzpicture}[scale=0.9]
          % Cable normal
          \draw[fill=lightbg, draw=darkcharcoal, line width=0.8pt] (0,0.6) rectangle (4.2,1.2);
          \draw[line width=1.8pt, utnblue] (0,0.95) -- (4.2,0.95) node[right, font=\tiny] {C1};
          \draw[line width=1.8pt, firered] (0,0.75) -- (4.2,0.75) node[right, font=\tiny] {C2};
          \node[fill=amberyellow!30, font=\fontsize{4}{5}\selectfont, inner sep=1pt] at (1.0,1.2) {Aislante Polímero};

          % Punto de calor y contacto
          \draw[fill=firered!30, draw=firered, line width=1pt, rounded corners=2pt] (2.3,0.5) rectangle (3.3,1.3);
          \draw[line width=2pt, darkcharcoal] (2.8,0.95) -- (2.8,0.75);
          \node[font=\fontsize{4.5}{5.5}\selectfont\bfseries, firered, above] at (2.8,1.3) {Calor: Fusión y Contacto};
          \node[font=\fontsize{4}{5}\selectfont\bfseries, utnblue, below] at (2.1,0.5) {$\Delta R \rightarrow$ Ubicación métrica exacta};
        \end{tikzpicture}
      }
    \end{column}

    \begin{column}{0.48\textwidth}
      \begin{block}{\textbf{Detectores Ópticos de Llama (UV / IR)}}
        \scriptsize
        $\bullet$ Sensores optoelectrónicos sintonizados al espectro infrarrojo y ultravioleta.\\
        $\bullet$ Algoritmos de parpadeo que discriminan la radiación solar y soldaduras.\\
        $\bullet$ Tiempo de respuesta \textbf{inferior a 1 segundo}.\\
        $\bullet$ \textbf{Uso:} Parques de tanques, turbinas y pintura.
      \end{block}
      \vspace{0.1cm}
      \centering
      \resizebox{0.95\columnwidth}{!}{%
        \begin{tikzpicture}[scale=0.9]
          % Detector en la pared
          \draw[fill=utnblue, draw=darkcharcoal, line width=1pt, rounded corners=2pt] (0,0.6) rectangle (0.8,1.4);
          \node[text=white, font=\fontsize{4.5}{5.5}\selectfont\bfseries, align=center] at (0.4,1.0) {Sensor\\UV/IR};

          % Cono de visión de detección
          \draw[fill=amberyellow!20, draw=fireorange, line width=0.8pt, opacity=0.7] (0.8,1.0) -- (3.8,1.8) -- (3.8,0.2) -- cycle;
          \node[font=\fontsize{4.5}{5.5}\selectfont\bfseries, fireorange] at (2.3,1.0) {Cono $90^\circ\text{--}120^\circ$};

          % Llama en el extremo
          \draw[fill=firered, draw=fireorange, line width=0.8pt] (3.9,0.7) .. controls (4.2,0.9) and (3.7,1.2) .. (4.0,1.4) .. controls (4.1,1.1) and (4.4,0.9) .. (3.9,0.7);
          \node[font=\fontsize{4.5}{5.5}\selectfont\bfseries, firered, right, align=left] at (4.1,1.0) {Llama\\$< 1\,\mathrm{s}$};
        \end{tikzpicture}
      }
    \end{column}
  \end{columns}
\end{frame}

% --- SLIDE 36: ANATOMÍA DEL EXTINTOR ---
\begin{frame}{Anatomía y Mecánica del Extintor Portátil}
  \begin{columns}[c]
    \begin{column}{0.54\textwidth}
      \small
      Componentes mecánicos (presión permanente):
      \begin{enumerate}
        \scriptsize
        \item \textbf{Cilindro metálico:} Chapa de acero para alta presión
        \item \textbf{Gas impulsor:} Nitrógeno seco ($N_2$) presurizado.
        \item \textbf{Tubo sifón:} Conduce el agente del fondo a la tobera.
        \item \textbf{Válvula de gatillo:} Apertura manual instantánea.
        \item \textbf{Manómetro indicador:} Aguja en zona verde ($\sim 14\,\mathrm{bar}$).
        \item \textbf{Pasador y precinto:} Seguro contra aperturas accidentales.
      \end{enumerate}
    \end{column}

    \begin{column}{0.44\textwidth}
      \centering
      \resizebox{0.65\columnwidth}{!}{%
        \begin{tikzpicture}
          % Cilindro
          \draw[fill=firered!85, draw=darkcharcoal, line width=1.5pt, rounded corners=5pt] (0,0) rectangle (2.4,3.5);
          % Tubo sifón interior
          \draw[dashed, line width=1.5pt, darkcharcoal!50] (1.2,0.2) -- (1.2,3.5);
          % Cabezal
          \draw[fill=darkcharcoal!80, line width=1.2pt] (0.8,3.5) rectangle (1.6,4.0);
          % Manija de transporte y palanca
          \draw[line width=2.5pt, darkcharcoal] (0.5,3.9) -- (1.8,4.3);
          \draw[line width=2.5pt, darkcharcoal] (0.5,3.7) -- (1.8,3.7);
          % Manómetro
          \draw[circle, fill=cardbg, draw=darkcharcoal, line width=1pt] (2.0,3.8) circle (0.24);
          \draw[line width=1pt, successgreen] (2.0,3.8) -- (2.0,4.0);
          % Manguera
          \draw[line width=2.5pt, darkcharcoal] (0.8,3.7) .. controls (-0.5,3.0) and (-0.4,1.8) .. (-0.1,1.0);
          % Etiqueta central
          \node[fill=white, draw=darkcharcoal, line width=0.8pt, font=\tiny\bfseries, align=center, inner sep=2pt] at (1.2,1.8) {Carga PQS /\\$\mathrm{CO_2}$};
        \end{tikzpicture}
      }
    \end{column}
  \end{columns}
\end{frame}

% --- SLIDE 37: PQS ABC Y BC ---
\begin{frame}{Extintores de Polvo Químico Seco (PQS ABC y BC)}
  \begin{columns}[T]
    \begin{column}{0.48\textwidth}
      \begin{block}{\textbf{PQS Clase ABC (Fosfato Monoamónico)}}
        \scriptsize
        $\bullet$ Base de $\mathrm{NH_4H_2PO_4}$ finamente micronizado.\\
        $\bullet$ \textbf{Fuegos Clase A:} Descompone a $>180\,\Celsius$ formando capa vidriada que sella poros y aísla el oxígeno.\\
        $\bullet$ \textbf{Fuegos Clase B:} Interrumpe reacción en cadena.\\
        $\bullet$ \textbf{Riesgo:} Altamente corrosivo en circuitería.
      \end{block}
      \vspace{0.1cm}
      \centering
      \resizebox{0.90\columnwidth}{!}{%
        \begin{tikzpicture}[scale=0.85]
          % Brasa sólida de madera
          \draw[fill=firered!70, draw=darkcharcoal, line width=1pt] (0,0) rectangle (2.8,0.7);
          \node[font=\fontsize{4.5}{5.5}\selectfont\bfseries, white] at (1.4,0.35) {Brasa Clase A};
          % Capa vidriada
          \draw[fill=utncyan!40, draw=utncyan!90!black, line width=1.2pt] (0,0.7) rectangle (2.8,1.0);
          \node[font=\fontsize{4.5}{5.5}\selectfont\bfseries, utncyan!90!black] at (1.4,0.85) {Capa Vidriada $\mathrm{HPO_3}$};
          % Advertencia corrosión
          \node[draw=firered, fill=firered!10, rounded corners=2pt, font=\fontsize{4}{5}\selectfont\bfseries, text=firered, align=center] at (1.4,-0.4) {$\triangle$ Residuo ácido daña plaquetas};
        \end{tikzpicture}
      }
    \end{column}

    \begin{column}{0.48\textwidth}
      \begin{block}{\textbf{PQS Clase BC (Bicarbonato)}}
        \scriptsize
        $\bullet$ Formulado con $\mathrm{NaHCO_3}$ o bicarbonato de potasio $\mathrm{KHCO_3}$.\\
        $\bullet$ \textbf{Mecanismo:} Captura química masiva de radicales libres en llama gaseosa sin fundirse.\\
        $\bullet$ \textbf{Ventaja:} Limpieza más fácil.
      \end{block}
      \vspace{0.1cm}
      \centering
      \resizebox{0.90\columnwidth}{!}{%
        \begin{tikzpicture}[scale=0.85]
          % Líquido inflamable Clase B
          \draw[fill=fireorange!30, draw=darkcharcoal, line width=1pt] (0,0) rectangle (2.8,0.6);
          \node[font=\fontsize{5}{6}\selectfont\bfseries, darkcharcoal] at (1.4,0.3) {Líquido Inflamable Clase B};
          % Llama y nube de polvo
          \draw[fill=fireorange!80, opacity=0.7] (0.4,0.6) .. controls (0.8,1.3) .. (1.4,0.9) .. controls (2.0,1.4) .. (2.4,0.6) -- cycle;
          \node[font=\fontsize{4.5}{5.5}\selectfont\bfseries, white] at (1.4,0.8) {Llama};
          \node[draw=successgreen!80!black, fill=successgreen!10, rounded corners=2pt, font=\fontsize{4}{5}\selectfont\bfseries, text=successgreen!80!black, align=center] at (1.4,-0.4) {Inhibición de radicales libres sin costra};
        \end{tikzpicture}
      }
    \end{column}
  \end{columns}
\end{frame}

% --- SLIDE 38: CO2 Y AGENTES LIMPIOS ---
\begin{frame}{Dióxido de Carbono ($CO_2$) y Agentes Limpios en Electrónica}
  \begin{columns}[c]
    \begin{column}{0.50\textwidth}
      Protección de circuitos energizados y centros de cómputo:
      \begin{itemize}
        \small
        \item \textbf{Dióxido de Carbono ($CO_2$):}
        \begin{itemize}
          \scriptsize
          \item Cilindros de acero sin costura a $50\text{--}60\,\mathrm{bar}$.
          \item Descarga en nieve carbónica a $\mathbf{-78.5\,\Celsius}$.
          \item Gas dieléctrico que se evapora al $100\%$ sin humedad ni residuos.
        \end{itemize}
        \item \textbf{Fluidos Fluorados:}
        \begin{itemize}
          \scriptsize
          \item Líquido vaporizante con cero potencial de daño a la capa de ozono.
          \item Enfriamiento molecular sin choque térmico extremo sobre silicio.
        \end{itemize}
      \end{itemize}
    \end{column}

    \begin{column}{0.48\textwidth}
      \centering
      \resizebox{0.92\columnwidth}{!}{%
        \begin{tikzpicture}[scale=0.9]
          % Tobera difusora cónica no conductora
          \draw[fill=darkcharcoal, draw=darkcharcoal, line width=1pt] (0.2,2.0) -- (1.6,2.4) -- (1.6,1.4) -- (0.2,1.8) -- cycle;
          \node[font=\fontsize{4}{5}\selectfont\bfseries, text=white] at (0.9,1.9) {Tobera Cónica};
          \draw[line width=3pt, darkcharcoal!70] (-0.3,1.9) -- (0.2,1.9);

          % Nube de CO2 / Gas limpio
          \draw[fill=utncyan!25, draw=utncyan, line width=0.8pt, opacity=0.85] (1.6,2.4) .. controls (2.4,2.8) and (3.4,2.7) .. (4.0,2.5) -- (4.0,1.2) .. controls (3.4,1.0) and (2.4,1.1) .. (1.6,1.4) -- cycle;
          \node[font=\fontsize{5}{6}\selectfont\bfseries, utncyan!90!black, align=center] at (2.8,1.9) {Nieve Carbónica\\$-78.5\,\Celsius$ / Gas};

          % Rack o componente electrónico
          \draw[fill=lightbg, draw=darkcharcoal, line width=1pt] (4.0,0.8) rectangle (4.6,2.8);
          \draw[line width=0.8pt, utnblue] (4.1,2.5) -- (4.5,2.5);
          \draw[line width=0.8pt, utnblue] (4.1,2.1) -- (4.5,2.1);
          \draw[line width=0.8pt, utnblue] (4.1,1.7) -- (4.5,1.7);
          \draw[line width=0.8pt, utnblue] (4.1,1.3) -- (4.5,1.3);
          \node[font=\fontsize{4}{5}\selectfont\bfseries, darkcharcoal, rotate=90] at (4.8,1.8) {Servidor};

          % Cuadro de propiedades
          \node[draw=successgreen!80!black, fill=successgreen!10, rounded corners=3pt, font=\fontsize{4.5}{5.5}\selectfont\bfseries, text=darkcharcoal, align=center] at (2.4,0.3) {Dieléctrico $\cdot$ Cero residuos $\cdot$ Sin daño electrónico};
        \end{tikzpicture}
      }
    \end{column}
  \end{columns}
\end{frame}

% --- SLIDE 39: POTENCIAL EXTINTOR ENSAYOS ---
\begin{frame}{Potencial Extintor y Ensayos de Homologación (IRAM 3517-2)}
  \begin{columns}[T]
    \begin{column}{0.48\textwidth}
      \begin{block}{\textbf{Rating Clase A ($1\mathrm{A}$ a $40\mathrm{A}$)}}
        \scriptsize
        $\bullet$ Ensayo sobre pira cúbica normalizada de listones de madera de pino secados en horno y dispuestos en capas cruzadas.\\
        $\bullet$ \textbf{Equivalencia 1A:} Extingue pira patrón de 50 listones equivalente a $5L$ de agua.\\
        $\bullet$ Un extintor $4\mathrm{A}$ apaga cuatro veces más masa de madera que uno de $1\mathrm{A}$.
      \end{block}
      \vspace{0.1cm}
      \centering
      \resizebox{0.85\columnwidth}{!}{%
        \begin{tikzpicture}[scale=0.8]
          % Pira de listones cruzados
          \foreach \y in {0, 0.25, 0.5, 0.75, 1.0} {
            \draw[fill=amberyellow!50, draw=darkcharcoal, line width=0.6pt] (0.2,\y) rectangle (2.6,\y+0.15);
          }
          \foreach \x in {0.4, 0.9, 1.4, 1.9, 2.4} {
            \draw[fill=amberyellow!80!black, draw=darkcharcoal, line width=0.6pt] (\x-0.1,0.05) rectangle (\x+0.1,1.1);
          }
          \node[font=\fontsize{5}{6}\selectfont\bfseries, firered, above] at (1.4,1.15) {Pira Normalizada Clase A};
          \node[draw=utnblue, fill=utnblue!10, rounded corners=2pt, font=\fontsize{4.5}{5.5}\selectfont\bfseries, text=utnblue] at (1.4,-0.4) {Rating: Capacidad de masa apagada};
        \end{tikzpicture}
      }
    \end{column}

    \begin{column}{0.48\textwidth}
      \begin{block}{\textbf{Rating Clase B ($1\mathrm{B}$ a $120\mathrm{B}$)}}
        \scriptsize
        $\bullet$ Ensayo sobre bateas cuadradas de acero con nafta en combustión libre.\\
        $\bullet$ \textbf{Definición del número:} Superficie de líquido inflamable que un operador es capaz de sofocar.\\
      \end{block}
      \vspace{0.1cm}
      \centering
      \resizebox{0.85\columnwidth}{!}{%
        \begin{tikzpicture}[scale=0.8]
          % Batea de acero
          \draw[fill=darkcharcoal!40, draw=darkcharcoal, line width=1pt] (0.2,0) -- (0.4,-0.3) -- (2.6,-0.3) -- (2.8,0) -- cycle;
          \draw[fill=fireorange!40, draw=fireorange, line width=0.8pt] (0.2,0) rectangle (2.8,0.25);
          % Llamas en batea
          \draw[fill=firered!80, opacity=0.8] (0.5,0.25) .. controls (0.9,1.1) .. (1.5,0.7) .. controls (2.1,1.2) .. (2.5,0.25) -- cycle;
          \node[font=\fontsize{5}{6}\selectfont\bfseries, firered, above] at (1.5,0.95) {Batea Líquida ($n$-heptano)};
          \node[draw=fireorange, fill=fireorange!10, rounded corners=2pt, font=\fontsize{4.5}{5.5}\selectfont\bfseries, text=fireorange!90!black] at (1.5,-0.65) {$10\mathrm{B} = 0.9\,\m^2$ sofocado por no experto};
        \end{tikzpicture}
      }
    \end{column}
  \end{columns}
\end{frame}

% --- SLIDE 40: RED DE HIDRANTES ---
\begin{frame}{Sistemas Fijos: Red Interior de Hidrantes}
  \begin{columns}[c]
    \begin{column}{0.50\textwidth}
      Instalación fija de suministro hídrico presurizado:
      \begin{itemize}
        \small
        \item \textbf{Reserva Exclusiva de Incendio:} Cisterna calculada para 60 a 120 minutos de agua.
        \item \textbf{Sala de Bombas Presurizadora:}
        \begin{itemize}
          \scriptsize
          \item Bomba principal eléctrica + Bomba de respaldo diésel.
          \item Bomba Jockey auxiliar para mantener la presión estática de red.
        \end{itemize}
        \item \textbf{Columna montante y gabinetes:}
        \begin{itemize}
          \scriptsize
          \item Mangueras de fibra sintética o devanaderas.
          \item Lanzas de chorro pleno y cono de niebla.
        \end{itemize}
        \item \textbf{Boca de impulsión exterior:} Para acople de autobomba de bomberos.
      \end{itemize}
    \end{column}

    \begin{column}{0.48\textwidth}
      \centering
      \resizebox{0.95\columnwidth}{!}{%
        \begin{tikzpicture}[scale=0.85]
          % Cisterna REI
          \draw[fill=utncyan!25, draw=utnblue, line width=1.2pt, rounded corners=2pt] (0,0) rectangle (1.7,1.6);
          \node[font=\fontsize{4.5}{5.5}\selectfont\bfseries, utnblue, align=center] at (0.85,0.8) {Cisterna REI\\(Agua Exclusiva)};
          \draw[line width=1.5pt, utnblue] (1.7,0.5) -- (2.2,0.5);

          % Sala de Bombas
          \draw[fill=cardbg, draw=darkcharcoal, line width=1pt, rounded corners=2pt] (2.2,0.1) rectangle (4.2,2.1);
          \node[font=\fontsize{4.5}{5.5}\selectfont\bfseries, darkcharcoal] at (3.2,1.85) {Sala de Bombas};
          
          % Bomba eléctrica
          \draw[fill=firered!20, draw=firered, line width=0.8pt, rounded corners=1pt] (2.35,1.05) rectangle (4.05,1.55);
          \node[font=\fontsize{4}{5}\selectfont\bfseries, firered] at (3.2,1.3) {Bomba Principal};
          
          % Bomba jockey
          \draw[fill=amberyellow!20, draw=amberyellow!80!black, line width=0.8pt, rounded corners=1pt] (2.35,0.35) rectangle (4.05,0.85);
          \node[font=\fontsize{4}{5}\selectfont\bfseries, amberyellow!80!black] at (3.2,0.6) {Bomba Jockey};

          % Colector y Columna montante
          \draw[line width=2pt, utnblue, ->, >=stealth] (4.05,1.3) -- (4.6,1.3) -- (4.6,3.3) node[above, font=\fontsize{4.5}{5.5}\selectfont\bfseries, utnblue] {Montante};

          % Gabinete R.I.A. en planta
          \draw[fill=firered!15, draw=firered, line width=1pt, rounded corners=2pt] (4.6,2.4) -- (5.0,2.4) -- (5.0,1.8) -- (6.0,1.8) -- (6.0,2.8) -- (5.0,2.8) -- (5.0,2.4);
          \draw[fill=firered, draw=darkcharcoal, line width=0.8pt] (5.2,2.0) rectangle (5.8,2.6);
          \node[font=\fontsize{4}{5}\selectfont\bfseries, text=white] at (5.5,2.3) {R.I.A.};
          \node[font=\fontsize{4}{5}\selectfont, darkcharcoal, right] at (5.0,1.6) {Manguera $45\mm$};

          % Boca de impulsión
          \draw[line width=1.5pt, firered, <-, >=stealth] (4.6,0.7) -- (5.5,0.7) node[right, font=\fontsize{4}{5}\selectfont\bfseries, firered] {Siamesa Bomberos};
        \end{tikzpicture}
      }
    \end{column}
  \end{columns}
\end{frame}

% --- SLIDE 41: SPRINKLERS E INUNDACIÓN GASEOSA ---
\begin{frame}{Rociadores Automáticos e Inundación Gaseosa}
  \begin{columns}[T]
    \begin{column}{0.48\textwidth}
      \begin{block}{\textbf{Rociadores Automáticos}}
        \scriptsize
        $\bullet$ Cañerías presurizadas permanentemente bajo techo.\\
        $\bullet$ \textbf{Ampolla de cuarzo:} Expansión por calor del líquido interno que estalla a temperatura de disparo.\\
        $\bullet$ \textbf{Puntual:} Abre solo el cabezal sobre el foco.
      \end{block}
      \centering
      \resizebox{0.65\columnwidth}{!}{%
        \begin{tikzpicture}[scale=0.75]
          % Rosca superior
          \draw[fill=darkcharcoal!40, draw=darkcharcoal, line width=0.8pt] (0.7,2.2) rectangle (1.7,2.6);
          \node[font=\fontsize{4}{5}\selectfont\bfseries] at (1.2,2.4) {Rosca};
          % Brazos del rociador
          \draw[line width=1.5pt, darkcharcoal] (0.8,2.2) -- (0.6,1.0) -- (1.2,0.6) -- (1.8,1.0) -- (1.6,2.2);
          % Ampolla termosensible de cuarzo (rojo 68°C)
          \draw[fill=firered, draw=darkcharcoal, line width=0.8pt, rounded corners=2pt] (1.05,0.7) rectangle (1.35,1.9);
          \fill[white] (1.15,1.3) circle (0.06); % burbuja
          % Deflector inferior dentado
          \draw[fill=amberyellow!80!black, line width=1pt] (0.3,0.5) -- (2.1,0.5) -- (2.0,0.35) -- (0.4,0.35) -- cycle;
          \node[font=\fontsize{4}{5}\selectfont\bfseries, darkcharcoal, below] at (1.2,0.35) {Deflector Distribuidor};
        \end{tikzpicture}
      }
    \end{column}

    \begin{column}{0.48\textwidth}
      \begin{block}{\textbf{Inundación Total por Gas Limpio}}
        \scriptsize
        $\bullet$ Cilindros presurizados vinculados a toberas en recinto cerrado.\\
        $\bullet$ Activación cruzada por 2 zonas para evitar falsos disparos.\\
        $\bullet$ Descarga integral en $<10\,\mathrm{s}$ sin daño informático.
      \end{block}
      \centering
      \resizebox{0.65\columnwidth}{!}{%
        \begin{tikzpicture}[scale=0.75]
          % Cilindros de gas
          \foreach \x in {0.4, 1.0, 1.6} {
            \draw[fill=utnblue!80, draw=darkcharcoal, line width=0.8pt, rounded corners=2pt] (\x,0.2) rectangle (\x+0.45,1.8);
            \draw[fill=darkcharcoal] (\x+0.12,1.8) rectangle (\x+0.33,2.0);
          }
          \node[font=\fontsize{4.5}{5.5}\selectfont\bfseries, utnblue] at (1.2,0.0) {Batería de Gas};
          % Colector de descarga
          \draw[line width=2pt, darkcharcoal] (0.6,2.0) -- (2.4,2.0) -- (2.4,2.6) -- (3.6,2.6);
          % Tobera en techo
          \draw[fill=darkcharcoal] (3.4,2.4) -- (3.8,2.4) -- (3.9,2.2) -- (3.3,2.2) -- cycle;
          \draw[->, >=stealth, line width=1pt, utncyan, dashed] (3.6,2.2) -- (3.1,1.2);
          \draw[->, >=stealth, line width=1pt, utncyan, dashed] (3.6,2.2) -- (3.6,1.0);
          \draw[->, >=stealth, line width=1pt, utncyan, dashed] (3.6,2.2) -- (4.1,1.2);
          \node[font=\fontsize{4.5}{5.5}\selectfont\bfseries, utncyan!90!black, below] at (3.6,0.9) {Inundación $<10\,\mathrm{s}$};
        \end{tikzpicture}
      }
    \end{column}
  \end{columns}
\end{frame}
